% !TeX root = main.tex
\newpage
\section{ForceDynamics Source Flow}

\newcommand{\ForceFunctionSection}[1]{%
  \subsection{ForceDynamics:#1}%
  \label{force:#1}%
}
\newcommand{\ForceFunctionEntry}[2]{%
  \subsubsection{\texttt{\detokenize{#1}}}%
  \label{force:#1}%
  #2%
}

\subsection{High-Level Model}

\texttt{ForceDynamics} is a source-owned, current-state contact model.  Each
mobile source reads an immutable frame-start snapshot, builds its complete
particle, piston, and parametric-wall contact set, accumulates one local force,
and writes only its own velocity and position.  Physical radii are used
unconditionally.  Initial particle or wall overlap is a configuration error;
contact history never changes geometry or makes overlapping objects
transparent.

\begin{figure}[htbp]
\centering
\resizebox{0.98\textwidth}{!}{\begin{tikzpicture}[
    node distance=0.27in and 0.48in,
    function/.style={draw=black!60, fill=black!3, rounded corners=2pt,
        minimum width=2.18in, minimum height=0.28in, align=center,
        font=\ttfamily\scriptsize},
    entry/.style={function, fill=green!8, draw=green!40!black,
        line width=0.65pt},
    wall/.style={function, fill=violet!7, draw=violet!55!black},
    constraint/.style={function, fill=orange!10, draw=orange!60!black},
    force/.style={function, fill=blue!7, draw=blue!45!black},
    flow/.style={-{Stealth[length=5pt]}, draw=black!65, line width=0.45pt},
    call/.style={-{Stealth[length=4pt]}, draw=blue!55!black, dashed,
        line width=0.4pt}
]

\newcommand{\ForceFlowNode}[2]{\hyperref[force:#1]{\strut #2}}

\node[entry] (main) {per-source compute invocation};

\node[wall, below left=0.46in and 1.42in of main] (piston)
    {\ForceFlowNode{ProcessPistonCollision}{ProcessPistonCollision}};
\node[function, below=0.46in of main] (physical)
    {\ForceFlowNode{GetPhysicalParticleContact}{GetPhysicalParticleContact}};
\node[wall, below right=0.46in and 1.42in of main] (marker)
    {\ForceFlowNode{ParametricBoundaryMarkerApplies}{ParametricBoundaryMarkerApplies}};

\node[wall, below=of piston] (piston-eval)
    {\ForceFlowNode{EvaluatePistonWall}{EvaluatePistonWall}};
\node[function, below=of physical] (particle)
    {\ForceFlowNode{ProcessParticleCollision}{ProcessParticleCollision}};
\node[wall, below=of marker] (wall-contacts)
    {\ForceFlowNode{EvaluateParametricWallContacts}{EvaluateParametricWallContacts}};

\node[function, below=of particle] (particle-state)
    {\ForceFlowNode{InitializeContactState}{InitializeContactState}};
\node[wall, below=of wall-contacts] (wall-eval)
    {\ForceFlowNode{EvaluateParametricWallSegment}{EvaluateParametricWallSegment}};
\node[wall, below=of wall-eval] (wall-process)
    {\ForceFlowNode{ProcessParametricWallCollision}{ProcessParametricWallCollision}};
\node[wall, below=of wall-process] (wall-state)
    {\ForceFlowNode{InitializeParametricWallContactState}{InitializeParametricWallContactState}};

\node[force, below=2.00in of particle-state] (force)
    {\ForceFlowNode{AccumulateContactForce}{AccumulateContactForce}};
\node[entry, below=of force] (velocity)
    {\ForceFlowNode{CalcVelocity}{CalcVelocity}};
\node[constraint, below=of velocity] (maximum)
    {\ForceFlowNode{ApplySourceMaximumDepth}{ApplySourceMaximumDepth}};
\node[constraint, below=of maximum] (project)
    {\ForceFlowNode{ProjectSourceVelocityToContactSet}{ProjectSourceVelocityToContactSet}};
\node[entry, below=of project] (position)
    {\ForceFlowNode{CalcPosition}{CalcPosition}};

\draw[flow] (main) -- (piston);
\draw[flow] (main) -- (physical);
\draw[flow] (main) -- (marker);

\draw[flow] (piston) -- (piston-eval);
\draw[flow] (physical) -- (particle);
\draw[flow] (particle) -- (particle-state);
\draw[flow] (marker) -- (wall-contacts);
\draw[flow] (wall-contacts) -- (wall-eval);
\draw[flow] (wall-eval) -- (wall-process);
\draw[flow] (wall-process) -- (wall-state);

\draw[call] (piston-eval.south) |- (force.west);
\draw[call] (particle-state.south) -- (force.north);
\draw[call] (wall-state.south) |- (force.east);
\draw[flow] (force) -- (velocity);
\draw[flow] (velocity) -- (maximum);
\draw[flow] (maximum) -- (project);
\draw[flow] (project) -- (position);

\end{tikzpicture}
}
\caption{Clickable shader-shaped flow exported from \texttt{ForceDynamics.py}.}
\label{fig:force-dynamics-flow}
\end{figure}

The diagram intentionally omits the Python harness, configuration loading,
scenario hooks, and diagnostic passes.  It shows the per-source function flow
that the exporter mirrors in GLSL.  Those Python-only support functions remain
documented later in this section but are not shader stages.

The stationary-wall representation is a set of finite cubic parametric
segments from \texttt{curve\_wall\_segments}.  Straight walls use zero
quadratic and cubic coefficients.  Boundary particles have positive
\texttt{ptype} and are locality markers only; no wall evaluator is carried in
\texttt{Data.z}.  Once any marker is local, the analytic segments supply the
wall point, tangent, normal, distance, penetration, and physical wall flag.

\clearpage
\subsection{Frame Orchestration}

\ForceFunctionEntry{CollisionRun}{Runs one complete semi-implicit frame.  It
gates the frame, executes scenario hooks, resets transient state, builds the
snapshot, validates initial geometry once, detects contacts, resolves velocity
and maximum depth, integrates position, records optional diagnostics, and
swaps the active buffers.}

\ForceFunctionEntry{BeginFrame}{Clears per-frame piston and error reporting.
A persistent configuration error stops the frame before contact processing.}

\ForceFunctionEntry{ApplyBeforeContactScanHook}{Calls the optional scenario
\texttt{BeforeContactScan} callback before snapshot construction.}

\ForceFunctionEntry{ResetFrameState}{Clears current-frame contact slots,
overlap measurements, and report fields while preserving recoverable internal
momentum and other cross-frame contact memory.}

\ForceFunctionEntry{ApplyStartRunHook}{Calls the optional scenario start hook
used by focused Python test cases.}

\ForceFunctionEntry{BuildFrameSnapshot}{Copies current positions and
velocities into immutable frame-start arrays used by all source calculations.}

\ForceFunctionEntry{BuildEndPositionSnapshot}{Captures integrated positions
only when diagnostics are enabled.  It does not trigger another force scan.}

\ForceFunctionEntry{DetectContacts}{Selects the active contact-determination
implementation.  The current Python model delegates to
\texttt{NaiveContactDetermination}.}

\ForceFunctionEntry{NaiveContactDetermination}{For every active mobile source,
processes the piston, scans mobile and boundary-marker targets, processes
physical particle contacts, then evaluates the deepest parametric contact for
each local physical wall flag.}

\ForceFunctionEntry{CalculateVelocities}{Calls \texttt{CalcVelocity} once for
each active mobile source after that source's complete force sum is known.}

\ForceFunctionEntry{CalculatePositions}{Calls \texttt{CalcPosition} once for
each active mobile source using the newly calculated velocity.}

\ForceFunctionEntry{EndFrame}{Toggles the active position/velocity buffer and
invalidates the Python frame snapshots.}

\subsection{Particle Contacts}

\ForceFunctionEntry{GetPhysicalParticleContact}{Returns validated geometry
only when two real-radius particles physically overlap.  It also detects a
source perimeter that has passed the target center and reports particle
tunneling.}

\ForceFunctionEntry{GetParticleGeometry}{Calculates center distance, contact
normal, overlap area, penetration inputs, and both physical radii.  No
effective or reduced radius exists.}

\ForceFunctionEntry{ProcessParticleCollision}{Creates or reuses the
source-owned contact state from previously calculated geometry and accumulates
its force without recalculating geometry.}

\ForceFunctionEntry{InitializeContactState}{Initializes the current-frame
particle contact, marks the source as colliding, and records contact
parameters.  The derived implementation uses the base contact ledger.}

\ForceFunctionEntry{GetContactState}{Finds an existing source-target contact
slot or appends one, then updates it with current physical geometry.}

\ForceFunctionEntry{AppendContactSlot}{Allocates a source-owned particle
contact slot, enforcing the configured slot capacity.}

\ForceFunctionEntry{ParticleContactTargetIDs}{Derives current particle target
IDs directly from active contact slots; there is no separate collision list.}

\ForceFunctionEntry{ParticlePenetrationDepth}{Returns
\(\max(0,r_s+r_t-d)\) for physical radii and center distance.}

\ForceFunctionEntry{particle_overlap_area}{Computes the exact circular lens
area for two intersecting particles, including containment and separation
limits.}

\subsection{Parametric Walls}

\ForceFunctionEntry{ParametricBoundaryMarkerApplies}{Checks boundary-marker
\texttt{ptype} and cell-local proximity.  Markers gate wall evaluation but do
not define wall geometry or select an evaluator.}

\ForceFunctionEntry{EvaluateParametricWallSegment}{Finds the closest bounded
point on one cubic segment, derives its tangent and outward normal, and returns
wall-contact geometry when the source overlaps the equal-radius ghost wall.}

\ForceFunctionEntry{EvaluateParametricWallContacts}{Evaluates all configured
segments and retains only the deepest valid contact for each positive physical
\texttt{wall\_flag}.}

\ForceFunctionEntry{ProcessParametricWallCollision}{Initializes one selected
wall contact and sends it through the common force path.}

\ForceFunctionEntry{InitializeParametricWallContactState}{Checks timestep
resolution and wall maximum depth, converts wall geometry into a source-owned
contact slot, and supplies stationary target velocity.}

\ForceFunctionEntry{AppendWallContactSlot}{Allocates or reuses a contact slot
keyed by physical wall flag rather than boundary-marker particle ID.}

\ForceFunctionEntry{ParametricWallPhysicalPenetration}{Calculates real-radius
penetration against a segment for initial-geometry validation.}

\ForceFunctionEntry{WallContactOffsetDistance}{Returns the configured radial
offset used to place the analytic equal-radius wall ghost.}

\ForceFunctionEntry{CheckWallMaximumDepth}{Reports wall tunneling when the wall
ghost's leading edge has passed the source center.}

\subsection{Piston}

\ForceFunctionEntry{PistonEnabled}{Returns true only when chamber bounds,
piston velocity, and piston start frame are all configured.}

\ForceFunctionEntry{GetPistonPosition}{Returns the piston plane position.  It
is stationary before its start frame, advances with timestep and configured
velocity, and clamps at the chamber end.}

\ForceFunctionEntry{GetPistonVelocity}{Returns zero while parked or clamped and
the configured velocity while the piston is moving.}

\ForceFunctionEntry{EvaluatePistonWall}{Evaluates the moving vertical ghost
plane within the chamber cross-section and returns common wall geometry.}

\ForceFunctionEntry{ProcessPistonCollision}{Checks the moving-wall timestep and
depth limits, creates a wall slot with piston target velocity, and accumulates
force and piston momentum diagnostics.}

\subsection{Contact Force and Memory}

\ForceFunctionEntry{AccumulateContactForce}{Applies stiffness along the contact
normal, updates compression or release phase, and stores recoverable inward
momentum on the source-owned contact.  The base implementation also records
detailed diagnostics.}

\ForceFunctionEntry{GetPairStiffness}{Returns the nonnegative mean of the two
mobile particles' stiffness values.}

\ForceFunctionEntry{GetContactStiffness}{Selects pair stiffness for particle
contacts and source stiffness for wall contacts.}

\ForceFunctionEntry{GetStartFrameVelocity}{Returns velocity from the immutable
frame snapshot, preventing source-order dependence.}

\ForceFunctionEntry{RecordRecoverableInternalMomentum}{Stores unresolved
inward normal momentum on the owning contact slot.}

\ForceFunctionEntry{SyncInternalMomentumMagnitude}{Recomputes the particle's
reported internal-momentum magnitude from its persistent contact slots.}

\ForceFunctionEntry{ClearRecoverableInternalMomentum}{Clears all recoverable
momentum owned by one particle.}

\ForceFunctionEntry{ClearInactiveRecoverableInternalMomentum}{Clears momentum
whose owning contact did not remain active in the current frame.}

\subsection{Maximum Depth and Timestep Safety}

The maximum allowed penetration is
\[
 d_{\max}=f_{\max}r_s,
\]
where the current maximum fraction is \(f_{\max}=0.5\).  A frame must also
satisfy
\[
 \Delta x_{\mathrm{inward}} \le (1-f_{\max})r_s.
\]

\ForceFunctionEntry{CheckPenetrationStepResolution}{Checks the relative inward
normal displacement for one timestep against the unused penetration reserve.
Failure reports \texttt{ERROR\_PENETRATION\_STEP\_TOO\_LARGE}.}

\ForceFunctionEntry{ApplySourceMaximumDepth}{Builds all currently active
particle and wall velocity constraints at maximum depth and applies them to
the source's candidate velocity as one source-owned operation.}

\ForceFunctionEntry{ProjectSourceVelocityToContactSet}{Projects the candidate
velocity onto at most two independent 2D normal constraints.  It merges nearly
parallel limits, rejects non-finite or ill-conditioned solutions, and returns
failure when no stable projection exists.  Its local \texttt{dot} and
\texttt{satisfies} helpers implement vector testing within this solver.}

\ForceFunctionEntry{SolveSmallLinearSystem}{Solves the bounded one- or two-row
constraint system with explicit singularity and finite-value checks.}

\subsection{Velocity, Position, and Activity}

\ForceFunctionEntry{CalcVelocity}{Performs semi-implicit Euler velocity update
from the complete source force, applies maximum-depth containment, and records
net source impulse diagnostics.}

\ForceFunctionEntry{CalcPosition}{Writes the alternate position buffer using
the newly calculated velocity, checks cell-array bounds, and marks mobile
particles dead when they leave the six death planes.}

\ForceFunctionEntry{GetParticleMass}{Returns \texttt{parms.x} with an epsilon
lower bound.}

\ForceFunctionEntry{GetParticlePosition}{Returns the immutable frame-start
position when available, otherwise the currently selected position buffer.}

\ForceFunctionEntry{GetCurrentParticlePosition}{Returns directly from the
currently selected position buffer without consulting the diagnostic snapshot.}

\ForceFunctionEntry{GetParticleVelocity}{Returns the velocity associated with
the active position buffer.}

\ForceFunctionEntry{particle_position}{Static buffered-position accessor used
before a full dynamics object snapshot is available.}

\ForceFunctionEntry{VelocityAngle}{Returns the planar velocity heading used by
rendering and reports.}

\ForceFunctionEntry{IsNullParticle}{Identifies reserved particle index zero.}

\ForceFunctionEntry{IsBoundaryParticle}{Identifies evaluator-bearing boundary
markers.}

\ForceFunctionEntry{IsMobileParticleActiveForDynamics}{Accepts only non-null,
non-boundary particles whose state is not dead.}

\subsection{Initial Validation and Errors}

\ForceFunctionEntry{ValidateInitialGeometry}{Runs once after configuration.
It rejects initial mobile-mobile overlap, parametric-wall overlap, and piston
overlap using physical geometry.}

\ForceFunctionEntry{SetError}{Stores the error code, category, source ID,
target ID, and wall flag in the shared collision result and returns false for
direct propagation through the call chain.}

\ForceFunctionEntry{ErrorDescription}{Formats the current error state for the
Python UI and report output.}

\subsection{Contact Ledger}

\ForceFunctionEntry{BeginContactFrame}{Starts a new frame for one source's
contact ledger and clears only transient slot activity.}

\ForceFunctionEntry{GetContactSlots}{Returns the fixed source-owned slot array.
The derived override exposes the same ledger to active dynamics code.}

\ForceFunctionEntry{ClearContactSlot}{Resets one reusable slot to inactive
defaults.}

\ForceFunctionEntry{ClearContactDiagnostics}{Clears diagnostic fields without
changing the physical identity of a persistent contact.}

\ForceFunctionEntry{RecordContactParameters}{Copies current geometry, phase,
velocity, force, and impulse values into report fields.}

\subsection{Diagnostics and Energy}

These functions run only when \texttt{dynamics\_diagnostics=true}; rendering
mode does not implicitly enable them.

\ForceFunctionEntry{ParticleKineticEnergy}{Returns one particle's kinetic
energy using either a supplied velocity or current buffered velocity.}

\ForceFunctionEntry{GetContactPotentialEnergy}{Returns elastic potential energy
for one particle contact under the configured force measure.}

\ForceFunctionEntry{GetOverlapPotentialEnergy}{Integrates the selected overlap
force measure to obtain contact potential energy.}

\ForceFunctionEntry{TotalPotentialEnergy}{Sums active particle-contact
potential energy without double counting physical pairs.}

\ForceFunctionEntry{TotalWallPotentialEnergy}{Sums active piston and parametric
wall potential energy.}

\ForceFunctionEntry{CurrentDiagnosticTotalEnergy}{Combines current kinetic,
particle-contact, and wall-contact energy.}

\ForceFunctionEntry{CapturePairGeometryDiagnostics}{Captures physical pair
distance, penetration, and phase information before or after integration.}

\ForceFunctionEntry{InitializePairPhaseEnergyReference}{Initializes the energy
reference used to measure one compression/release phase.}

\ForceFunctionEntry{RecordPairPhaseDiagnostics}{Compares start and end pair
geometry and records phase-level energy changes.}

\ForceFunctionEntry{RecordFrameStartDiagnostics}{Captures frame-start momentum,
kinetic energy, potential energy, and piston state.}

\ForceFunctionEntry{RecordAfterResolveDiagnostics}{Captures the corresponding
post-resolution quantities.}

\subsection{Configuration and Construction}

\ForceFunctionEntry{load_cfg_file}{Loads constants, the selected cfg and its
include, particle data, marker data, shader flags, dimensions, force mode,
diagnostic switches, and persistent configuration errors.}

\ForceFunctionEntry{load_include}{Merges the configured include file into the
primary configuration object.}

\ForceFunctionEntry{load_constants}{Loads named error codes and shared numeric
constants.}

\ForceFunctionEntry{add_parametric_boundary_markers_from_cfg}{Builds
evaluator-5 locality markers when particles are supplied directly by cfg.}

\ForceFunctionEntry{getParticleData}{Reads the selected binary particle file
and preserves its null particle and buffered fields.}

\ForceFunctionEntry{create_particle_from_cfg}{Converts one cfg particle
record into the runtime particle structure.}

\ForceFunctionEntry{create_particle_array_from_cfg}{Builds the indexed
runtime particle array from binary or cfg records.}

\ForceFunctionEntry{create_particle_array}{Allocates an empty particle array
with the requested count.}

\ForceFunctionEntry{create_particle}{Constructs one initialized particle
record with buffered position and velocity fields.}

\ForceFunctionEntry{create_shader_flags_from_cfg}{Builds frame, timestep,
buffer-selector, and rectangular cell-array state from configuration.}

\ForceFunctionEntry{create_shader_flags}{Constructs an initialized shader
flags object.}

\ForceFunctionEntry{create_collision_in}{Constructs the shared collision
result and error-reporting object.}

\ForceFunctionEntry{create_geo_contact_state}{Constructs one initialized
source-owned contact slot.}

\ForceFunctionEntry{create_vec4}{Constructs the Python four-component float
vector used to mirror GLSL structures.}

\ForceFunctionEntry{create_uvec4}{Constructs the corresponding unsigned
integer vector.}

The small structure classes at the beginning of \texttt{ForceDynamicsBase.py}
use \texttt{\_\_setattr\_\_} only to provide GLSL-like aliasing and field
normalization.  They are storage mechanics, not dynamics stages, and therefore
do not appear as flowchart nodes.
